\documentclass[12pt]{article}
\providecommand{\C}{\mathbb{C}}
\providecommand{\Z}{\mathbb{Z}}
\providecommand{\H}{\mathbb{H}}
\providecommand{\gDelta}{\mathfrak{g}_{\Delta_5}}
\providecommand{\slhat}[1]{\widehat{\mathfrak{sl}}_{#1}}
\providecommand{\sl}[1]{\mathfrak{sl}_{#1}}
\providecommand{\LambII}{\Lambda^{2,1}_{II}}
\providecommand{\FF}{F_3}
\providecommand{\kBKM}{\kappa_{\mathrm{BKM}}}
\providecommand{\kch}{\kappa_{\mathrm{ch}}}
\begin{document}

\section*{Wave-15 M5/20: the $\eta^9$ exponent and $\slhat{3}$}

\subsection*{Claim under test}

$\mathfrak{g}_{\Delta_5}$ contains a level-$1$ affine subalgebra
$\slhat{3} \hookrightarrow \gDelta$ whose Weyl--Kac denominator
produces the prefactor $\eta^9$ in the zero-cusp factorisation
$A(z_1,z_2) = \eta(z_1)^9\,\nu_{11}(z_1,z_2)$ of
$\Delta_5 = \mathrm{Grit}(\eta^9 \vartheta_1)$ at $\kBKM = 5$.

\subsection*{Attack 1: Cartan matrix comparison}

The Kac--Moody Cartan of $\sl{3}$ is $A_{\sl{3}} =
\begin{pmatrix} 2 & -1 \\ -1 & 2 \end{pmatrix}$ (rank $2$,
positive definite, $\det = 3$). The
Gritsenko--Nikulin real simple roots $\{\delta_1,\delta_2,\delta_3\}$
of $\gDelta$ carry the Gram matrix
\[
  G = \begin{pmatrix} 2 & -2 & -2 \\ -2 & 2 & -2 \\ -2 & -2 & 2 \end{pmatrix},
  \qquad \det G = -32,
\]
with eigenvalues $\{+4,+4,-2\}$, signature $(2,1)$
(Feingold--Frenkel 1983 \emph{Math.\ Ann.}~263; Gritsenko--Nikulin
1998 \texttt{alg-geom/9712020} \S3). Rank and signature
disagree: $A_{\sl{3}}$ is rank $2$ positive definite; $G$ is
rank $3$ hyperbolic.

\emph{Heal.} The mismatch is structural, not notational. No
orthogonal transformation sends $G$ to $A_{\sl{3}} \oplus
(\text{null})$: a $1$-dimensional null subspace of $G$ would
require $\det G = 0$, contradicting $\det G = -32$.

\subsection*{Attack 2: counting roots}

$\sl{3}$ has $9 = 6 + 3$: six $\pm$ roots
$\{\pm\alpha_1,\pm\alpha_2,\pm(\alpha_1+\alpha_2)\}$ and three
Cartan generators. The exponent $9$ in $\eta^9$ is a numerical
match.

\emph{Heal.} $\eta^9 = \eta^1 \cdot \eta^8$ in the factorisation
at the zero-dimensional cusp (slide~803). The $\eta^1$ is the
Weyl vector prefactor $e^{\pi i(z_1+z_2+z_3)}$ evaluated along
the null direction $a_0 = \delta_1$ at pairing $(\rho,\delta_1)
= \tfrac{1}{2}$ (wave9 P2 \S2). The $\eta^8$ comes from the
Borcherds product
\[
  \prod_{n \geq 1}(1 - q_1^n)^{\sum_l f(0,l)}
  = \prod_{n \geq 1}(1 - q_1^n)^{10 - 1 - 1} = \eta^8 \cdot q_1^{-1/3},
\]
using Gritsenko--Nikulin Fourier coefficients $f(0,0) = 10$,
$f(0,\pm 1) = 1$, $f(0,|l|\geq 2) = 0$. The split $9 = 1 + 8$
therefore encodes one Weyl contribution plus eight transverse
oscillators, not the $9 = 6 + 3$ root-plus-Cartan split of
$\sl{3}$. The coincidence of the total is genuine, the
mechanism is different.

\subsection*{Attack 3: Weyl--Kac denominator of $\slhat{3}$ at level $1$}

The Kac 1990 \emph{Infinite-dim.\ Lie algebras} \S 12.7
denominator identity at level $k$ for $\slhat{3}$ reads
\[
  \prod_{n \geq 1}(1-q^n)^2 \prod_{\alpha \in \Delta^+_{\sl{3}}}
    \prod_{n \geq 0}(1 - q^n e^{-\alpha})(1 - q^{n+1} e^{\alpha})
  = \sum_{w \in \widehat{W}} \varepsilon(w)\, e^{w \rho - \rho}.
\]
Specialising to the character at the $\sl{3}$ origin $e^{\alpha} \to 1$:
the left side becomes $\prod_{n \geq 1}(1-q^n)^{2 + 2 \cdot 6} =
\eta(q)^{14}\,q^{-14/24}$, exponent $14 = 2 \cdot \mathrm{rank}
+ 2 \cdot |\Delta^+| = 2 + 12$. The string function $c^0_0(\tau)$
of $\slhat{3}$ at level $1$ is $\eta(\tau)^{-2}$ (Kac--Peterson
1984 Theorem~4.9); the Frenkel--Kac 1980 construction realises
$\slhat{3}$ level $1$ on the $A_2$-root lattice VOA with vacuum
character $\eta(q)^{-2}\,\Theta_{A_2}(q,z)$, not $\eta^9$.

\emph{Heal.} The $\slhat{3}$ level-$1$ vacuum character is
$\eta^{-2}\Theta_{A_2}$. Multiplied by $\eta^{11}$ (the
Frenkel--Kac bosonic-ghost contribution for the two free bosons
in signature-$(2,0)$ of $A_2$) the denominator is $\eta^{11}$,
not $\eta^9$. The exponent mismatch is $2$. No specialisation of
the $\slhat{3}$ level-$1$ denominator produces $\eta^9$ along a
rank-$1$ direction.

\subsection*{Attack 4: real-root sector of $\gDelta$ as identified in Wave~7}

Wave~7 b5 (\texttt{wave7\_b5\_vacuum\_OPE\_depth\_5.tex} l.~29)
identifies the real-root subalgebra of the vertex algebra
$V_{\Delta_5}$: structure constants $f^{ij}{}_k$ vanish because
$\delta_i + \delta_j \notin \Delta^{\mathrm{re}}$ for $i \neq
j$ (direct enumeration of $\mathcal{P}_{II}$, the primitive
real-root polyhedron). The real-root sector is therefore the
\emph{abelian} rank-$3$ hyperbolic Heisenberg current algebra
$\widehat{\mathfrak{h}}_{\LambII}$ of signature $(2,1)$:
\[
  [J^{\delta_i}_n, J^{\delta_j}_m] = n\,(\delta_i,\delta_j)\,\delta_{n+m,0},
\]
with $(\delta_i,\delta_j) = G_{ij}$ the GN Gram. An abelian
rank-$3$ indefinite Heisenberg is not $\slhat{3}$: the latter
is non-abelian of rank $2$ and level $k = 1$ central.

\emph{Heal.} The "$\FF$" in
$\gDelta = \mathrm{Borch}(\FF, \phi_{0,1}^{K3})$
(Adjudication Ledger W14--W19, entry~32) is the
Feingold--Frenkel rank-$3$ hyperbolic Kac--Moody, not affine
$\sl{3}$. The real-root subalgebra of $\gDelta$ is $\FF$, and
its currents are abelian because $\delta_i + \delta_j$ are
imaginary (norm $\leq -2$) hence outside the real-root lattice.
The imaginary roots, carrying multiplicity $|c_{\phi_{0,1}}
(D(\alpha))|$, are where the BKM non-abelianness lives.

\subsection*{Attack 5: consequence, the refuted candidate theorem}

A putative theorem "$\slhat{3}_{k=1} \hookrightarrow \gDelta$
as the real-root closed subalgebra" would require:
(i) the real simple roots of $\gDelta$ to form a rank-$2$
positive-definite $A_2$ sub-root-system; (ii) the real-root
bracket to be non-abelian. Attack~1 refutes (i) by rank and
signature; Attack~4 refutes (ii) by direct OPE vanishing.

\emph{Heal.} The extracted true theorem is different and
weaker: there exists a rank-$2$ positive-definite subspace
$P \subset \LambII \otimes \R$ (spanned by the two
$+4$-eigenvectors of $G$), and the projection
$\pi_P(\delta_i)$ of the real simple roots onto $P$ gives
three vectors of equal squared norm $4$ pairwise angled at
$120^\circ$, realising an $A_2$ root diagram rescaled by
$\sqrt{2}$ from the standard $A_2$-lattice. The Heisenberg
$\widehat{\mathfrak{h}}_{P}$ of signature $(2,0)$ on this
plane carries a free-boson realisation of $\slhat{3}_1$ via
Frenkel--Kac on the rescaled $A_2$, but this $\slhat{3}_1$
does \emph{not} embed into $\gDelta$ as a closed subalgebra
of the real-root sector: the rescaling from $A_2$ to
$\sqrt{2} A_2$ multiplies the level by $2$ and the
corresponding lattice VOA is $V_{\sqrt{2}A_2}$, not
$V_{A_2}$.

Scheithauer 2004 \emph{Contemp.\ Math.}~343 (``Generalized
Kac--Moody algebras and some related topics'') constructs
BKMs from lattice VOAs with BRST-like screening; in that
framework an $\slhat{3}_1$ screening would need three
exponential currents $e^{\alpha_i}$ with $(\alpha_i,\alpha_j) =
A_{\sl{3}}$, conflicting with the $G$-Gram of $\gDelta$.

\subsection*{Verdict}

\textbf{The $\slhat{3}$-embedding conjecture is refuted.} The
$\eta^9$ prefactor has a different, first-principles origin:
one Weyl vector contribution ($\eta^1$) plus eight transverse
Heisenberg oscillators from $f(0,0) - f(0,1) - f(0,-1) = 8$
($\eta^8$). The three real simple roots span a
signature-$(2,1)$ hyperbolic Feingold--Frenkel Cartan, not the
positive-definite rank-$2$ $A_2$ Cartan. Kac--Peterson
string-function arithmetic gives $\slhat{3}_1$ vacuum
$\eta^{-2}\Theta_{A_2}$, incompatible with $\eta^9$.

\textbf{True residual.} The $+4$-eigenspace $P \subset
\LambII \otimes \R$ carries an $A_2$-like configuration of
projected real simple roots, hinting at a Langlands-dual
\emph{positive-definite shadow} of $\FF$. This shadow is a
candidate for the $\sl{3}$-motivic-cohomology statement in
the geometric Satake-at-level-$1$ picture, but is not a
subalgebra.

\subsection*{Primary sources verified}

Kac 1990 \emph{Infinite-dimensional Lie algebras}~\S 12.7,
denominator identity. Frenkel--Kac 1980 \emph{Invent.}~62,
free-field $A$-$D$-$E$ realisation at level $1$.
Gritsenko--Nikulin 1998 \texttt{alg-geom/9712020} \S3,
$\Delta_5$ real simple roots and Gram $G$. Feingold--Frenkel
1983 \emph{Math.\ Ann.}~263, rank-$3$ hyperbolic Kac--Moody
$F_3$ Cartan eigenvalues $\{+4,+4,-2\}$. Scheithauer 2004
\emph{Contemp.\ Math.}~343, BKM from lattice VOA screening.

\end{document}
